\section*{Chapter \ref{ch:g9:sqrt} --- Square Roots}

\begin{solution}{exo:g9:sqrt:1}
$\sqrt{64} = 8$; \quad $\sqrt{100} = 10$; \quad
$\sqrt{0.09} = 0.3$ (since $0.3^2 = 0.09$); \quad
$\sqrt{\frac{1}{25}} = \frac15$; \quad
$\left(\sqrt{13}\right)^2 = 13$; \quad
$\sqrt{(-4)^2} = \sqrt{16} = 4 = \abs{-4}$.
\end{solution}

\begin{solution}{exo:g9:sqrt:2}
$\sqrt{50} = \sqrt{25 \times 2} = 5\sqrt2$; \quad
$\sqrt{45} = \sqrt{9 \times 5} = 3\sqrt5$; \quad
$\sqrt{98} = \sqrt{49 \times 2} = 7\sqrt2$; \quad
$\sqrt{300} = \sqrt{100 \times 3} = 10\sqrt3$.
\end{solution}

\begin{solution}{exo:g9:sqrt:3}
$\sqrt2 \times \sqrt{18} = \sqrt{36} = 6$.

$\dfrac{\sqrt{75}}{\sqrt3} = \sqrt{\dfrac{75}{3}} = \sqrt{25} = 5$.

$\sqrt5 \times \sqrt{20} = \sqrt{100} = 10$.

$\dfrac{\sqrt8}{\sqrt{50}} = \sqrt{\dfrac{8}{50}} = \sqrt{\dfrac{4}{25}}
= \dfrac25$.
\end{solution}

\begin{solution}{exo:g9:sqrt:4}
$x^2 = 36$: $x = 6$ or $x = -6$.

$x^2 = 11$: $x = \sqrt{11}$ or $x = -\sqrt{11}$.

$x^2 + 4 = 0$ means $x^2 = -4 < 0$: no solution.

$2x^2 = 50$: $x^2 = 25$, so $x = 5$ or $x = -5$.
\end{solution}

\begin{solution}{exo:g9:sqrt:5}
$\sqrt{20} + \sqrt{45} = 2\sqrt5 + 3\sqrt5 = 5\sqrt5$.

$3\sqrt8 - \sqrt{18} + \sqrt2 = 3 \times 2\sqrt2 - 3\sqrt2 + \sqrt2
= (6 - 3 + 1)\sqrt2 = 4\sqrt2$.
\end{solution}

\begin{solution}{exo:g9:sqrt:6}
$\left(\sqrt3 + 1\right)^2 = 3 + 2\sqrt3 + 1 = 4 + 2\sqrt3$.

$\left(2\sqrt5 - 3\right)^2 = 4 \times 5 - 2 \times 3 \times 2\sqrt5 + 9
= 29 - 12\sqrt5$.

$\left(\sqrt6 + \sqrt2\right)\left(\sqrt6 - \sqrt2\right) = 6 - 2 = 4$.
\end{solution}

\begin{solution}{exo:g9:sqrt:7}
Side: $\sqrt{6400} = 80$ m. Diagonal (Pythagoras):
$\sqrt{80^2 + 80^2} = 80\sqrt2 \approx 113$ m.
\end{solution}

\begin{solution}{exo:g9:sqrt:8}
$\dfrac{1}{\sqrt2} = \dfrac{1 \times \sqrt2}{\sqrt2 \times \sqrt2}
= \dfrac{\sqrt2}{2}$.

$\dfrac{6}{\sqrt3} = \dfrac{6\sqrt3}{3} = 2\sqrt3$.
\quad
$\dfrac{10}{\sqrt5} = \dfrac{10\sqrt5}{5} = 2\sqrt5$.
\end{solution}

\begin{solution}{exo:g9:sqrt:9}
\emph{1. True}: this is the \omterm{def:g2:mult:def}{product} rule (\cref{thm:g9:sqrt:rules}).

\emph{2. False}: $\sqrt{9 + 16} = 5$ but $\sqrt9 + \sqrt{16} = 7$.

\emph{3. False} for negative $x$: $\sqrt{(-4)^2} = 4 \neq -4$. The
correct statement is $\sqrt{x^2} = \abs{x}$.

\emph{4. True}: $\left(3\sqrt2\right)^2 = 9 \times 2 = 18$.
\end{solution}

\begin{solution}{exo:g9:sqrt:10}
$\left(2 + \sqrt3\right)^2 = 4 + 4\sqrt3 + 3 = 7 + 4\sqrt3$. So
$x = \sqrt{7 + 4\sqrt3} = \sqrt{\left(2+\sqrt3\right)^2} = 2 + \sqrt3$
(a nonnegative number, so the root just removes the square).

Similarly $\left(\sqrt5 - 2\right)^2 = 5 - 4\sqrt5 + 4 = 9 - 4\sqrt5$,
and $\sqrt5 - 2 > 0$, so
$\sqrt{9 - 4\sqrt5} = \sqrt5 - 2$ (not $2 - \sqrt5$, which is
negative!).
\end{solution}

\begin{solution}{pb:g9:sqrt:1}
\textbf{1.} $1.4^2 = 1.96 < 2 < 2.25 = 1.5^2$;
$1.41^2 = 1.9881 < 2 < 2.0164 = 1.42^2$. Three decimals:
$1.414^2 = 1.999396 < 2 < 2.002225 = 1.415^2$, so
$1.414 < \sqrt2 < 1.415$.

\textbf{2.} Squaring $\sqrt2 = \frac ab$ gives
$2 = \frac{a^2}{b^2}$, hence $a^2 = 2b^2$: the number $a^2$ is
twice a whole number --- \emph{\omterm{def:g3:numbers:evenodd}{even}}.

\textbf{3.} If $a$ were \omterm{def:g3:numbers:evenodd}{odd}, $a^2$ would be \omterm{def:g3:numbers:evenodd}{odd}
(\cref{pb:g8:pythagoras:1}); since $a^2$ is \omterm{def:g3:numbers:evenodd}{even}, $a$ is \omterm{def:g3:numbers:evenodd}{even}:
$a = 2k$. Substituting: $4k^2 = 2b^2$, so $b^2 = 2k^2$ is \omterm{def:g3:numbers:evenodd}{even},
and by the same parity fact $b$ is \omterm{def:g3:numbers:evenodd}{even}.

\textbf{4.} Both $a$ and $b$ \omterm{def:g3:numbers:evenodd}{even} means both \omterm{def:g6:wholes:divisible}{divisible} by $2$
--- but $\frac ab$ was in lowest terms, sharing no common
\omterm{def:g9:arith:divisor}{divisor}. Contradiction: the assumed \omterm{def:g9:fractions:fraction}{fraction} cannot exist.
\emph{Theorem: $\sqrt2$ is irrational.}

\textbf{5.} The contradiction lives entirely in ``lowest
terms'': without it, ``$a$ and $b$ both \omterm{def:g3:numbers:evenodd}{even}'' contradicts
nothing. Assuming lowest terms is legitimate because every
\omterm{def:g9:fractions:fraction}{fraction} \emph{has} a lowest-terms form
(\cref{met:g9:arith:simplify}: divide out the GCD) --- if
$\sqrt2$ were any \omterm{def:g9:fractions:fraction}{fraction} at all, it would be a lowest-terms
one, and that one is impossible.

\textbf{6.} In prime factorizations, squares carry \omterm{def:g3:numbers:evenodd}{even}
\omterm{def:g8:powers:def}{exponents} (\cref{pb:g9:arith:1}). In $a^2 = 3b^2$, the \omterm{def:g8:powers:def}{exponent}
of $3$ is \omterm{def:g3:numbers:evenodd}{even} on the left, but \omterm{def:g3:numbers:evenodd}{odd} on the right (\omterm{def:g3:numbers:evenodd}{even} from
$b^2$, plus one). No number has two factorizations
(\cref{thm:g9:arith:factorization}): contradiction, and
$\sqrt3$ is irrational.

\textbf{7.} In $a^2 = n b^2$, the argument finds a prime with
contradictory \omterm{def:g8:powers:def}{exponent} parity exactly when some prime appears in
$n$ with an \emph{\omterm{def:g3:numbers:evenodd}{odd}} \omterm{def:g8:powers:def}{exponent}. If every \omterm{def:g8:powers:def}{exponent} of $n$ is
\omterm{def:g3:numbers:evenodd}{even}, $n$ is a perfect square ($n = m^2$, and
$\sqrt n = m$ is whole). Complete result: $\sqrt n$ is
irrational for every whole $n$ that is not a perfect square ---
$\sqrt2, \sqrt3, \sqrt5, \sqrt6, \sqrt7, \sqrt8, \sqrt{10},
\dots$

\textbf{8.} From $r x = q$ with $r \neq 0$:
$x = \frac qr$, a \omterm{def:g3:division:remainder}{quotient} of rationals, hence rational. So if
$x$ is irrational, $r x$ cannot be rational: $2\sqrt2$ and
$\frac{\sqrt2}{2}$ are irrational (multipliers $2$ and
$\frac12$).

\textbf{9.} If $1 + \sqrt2 = q$ were rational, then
$\sqrt2 = q - 1$ would be rational: contradiction. If
$s = \sqrt2 + \sqrt3$ were rational, then
$s^2 = 2 + 2\sqrt6 + 3 = 5 + 2\sqrt6$ would be too, giving
$\sqrt6 = \frac{s^2 - 5}{2}$ rational --- but $\sqrt6$ is
irrational (question 7). So $\sqrt2 + \sqrt3$ is irrational.

\textbf{10.} \omterm{def:g1:addition:def}{Sums}: $\sqrt2$ and $-\sqrt2$ (or $\sqrt2$ and
$1 - \sqrt2$) are each irrational, with rational \omterm{def:g1:addition:def}{sums} $0$ and
$1$. \omterm{def:g2:mult:def}{Products}: $\sqrt2 \times \sqrt2 = 2$. Irrationality can
evaporate in arithmetic --- hence the case-by-case proofs.

\textbf{11.} From $x = 1$: average of $1$ and $\frac21 = 2$:
$x_1 = \frac32$. Then average of $\frac32$ and
$\frac{2}{3/2} = \frac43$:
$x_2 = \frac12\left(\frac32 + \frac43\right)
= \frac12 \cdot \frac{17}{6} = \frac{17}{12} \approx 1.4167$.

\textbf{12.} $x_3 = \frac12\left(\frac{17}{12} +
\frac{24}{17}\right) = \frac12 \cdot
\frac{289 + 288}{204} = \frac{577}{408} = 1.4142156\ldots$
Against $\sqrt2 = 1.4142135\ldots$: five decimals correct after
three steps --- the recipe roughly doubles the correct decimals
each \omterm{def:g4:numbers:round}{round}.

\textbf{13.} If $x^2 > 2$ (guess too big), then $x > \sqrt2$,
and $\frac2x < \frac{2}{\sqrt2} = \sqrt2$: the companion side is
too small. If $x^2 < 2$, the inequalities reverse. Either way
$\sqrt2$ sits between $x$ and $\frac2x$, so their average ---
the next guess --- is closer than the worse of the two sides:
the rectangle of \omterm{def:g6:measure:area}{area} $2$ gets squarer at every step.

\textbf{14.} $9 - 8 = 1$; \ $289 - 288 = 1$; \
$332\,929 - 332\,928 = 1$. Every Heron guess $\frac ab$
satisfies $a^2 - 2b^2 = 1$: Part I proved that hitting $0$ is
impossible, and these \omterm{def:g9:fractions:fraction}{fractions} miss the impossible by exactly
one unit --- the closest whole numbers can ever come.

\textbf{15.} The number line must contain numbers beyond the
\omterm{def:g9:fractions:fraction}{fractions} --- lengths like $\sqrt2$, whose decimal writings run
forever without repeating: the \emph{real numbers}. Their
honest construction is a story for the High School volume and,
in full rigor, the university ones.
\end{solution}
